\documentclass[conference]{IEEEtran}
\usepackage[UTF8, heading=true]{ctex}
\usepackage{amsmath,amssymb,amsfonts}
\usepackage{graphicx}
\usepackage{booktabs}
\usepackage{multirow}
\usepackage{array}
\usepackage{float}
\usepackage{cite}
\usepackage{hyperref}
\usepackage{enumitem}
\hypersetup{colorlinks=true,linkcolor=blue,citecolor=blue,urlcolor=blue}
\renewcommand{\figurename}{图}
\renewcommand{\tablename}{表}

\begin{document}

\title{面向多智能体团队的技能萃取与赋予系统\\
\large{A Skill Extraction and Assignment System for Multi-Agent Teams}}

\author{\IEEEauthorblockN{AgentsGroup2026 研究团队}\IEEEauthorblockA{AgentsGroup2026 项目组\\Email: research@agentsgroup2026.dev}}
\maketitle

\begin{abstract}
大语言模型（LLM）驱动的多智能体团队在执行复杂任务时需要可复用的专业知识——技能（Skill）。人工编写技能无法规模化应对日益增长的任务多样性。本文提出AgentsGroup2026平台中的完整技能生命周期管理系统，涵盖从对话会话中自动萃取技能、向量索引与意图感知检索、到生态压力下技能演化与跨代遗传的全流程。系统包含三个核心子系统：（1）基于WebSocket实时通信的七阶段技能萃取流水线，支持时间轴UI的实时推送；（2）采用TF-IDF向量索引结合时间衰减与频率增强策略的智能检索器；（3）自然选择驱动的技能演化引擎，通过证据收集、LLM改进与A/B测试实现技能指令的持续优化。在AWS运维与构建系统两个领域、9次演化运行、45个智能体的实验中，系统成功实现了技能的自动化萃取、索引、检索、演化与团队赋予，验证了技能跨代遗传的有效性。

\smallskip
\noindent\textbf{关键词：}技能萃取，多智能体系统，LLM智能体，自然选择，技能演化，向量索引，WebSocket实时通信
\end{abstract}

\begin{IEEEkeywords}
skill extraction, multi-agent teams, LLM agents, natural selection, skill evolution, vector indexing, WebSocket real-time communication
\end{IEEEkeywords}

% ═══════════════════════════════════════════
\section{引言}
% ═══════════════════════════════════════════

大语言模型驱动的智能体系统在复杂任务推理与协作方面取得了显著进展~\cite{wu2023autogen,hong2023metagpt,park2023generative}。然而，智能体在执行领域特定任务时，需要依赖于可复用的专业知识——即"技能"（Skill）——这些技能通常以文本指令的形式编码了领域知识、操作流程和决策逻辑。Voyager~\cite{wang2023voyager}率先在Minecraft开放世界中展示了LLM智能体通过技能库积累和复用行为知识的能力，其不断增长的代码技能库成为智能体持续探索的基础设施。

然而，Voyager的技能库由智能体自身生成代码构成，技能创建过程与智能体行为紧密耦合，缺乏对多智能体团队场景的泛化支持。在实际的多智能体协作系统中，一个根本性挑战浮现：\textbf{如何从多智能体交互会话中自动萃取可复用的技能，将其索引为可检索的知识库，并在生态压力下持续演化与赋予给团队成员？}

当前多智能体框架如AutoGen~\cite{wu2023autogen}、MetaGPT~\cite{hong2023metagpt}和Generative Agents~\cite{park2023generative}主要关注智能体间的对话编排与角色分工，但均缺乏系统性的技能生命周期管理机制。Reflexion~\cite{shinn2023reflexion}和ReAct~\cite{yao2022react}提供了智能体自我反思与推理的能力，但反思结果并未被结构化为可跨会话复用的技能。CoALA~\cite{sumers2023coala}从认知架构角度提出了语言智能体的理论框架，将技能作为长期记忆模块，但未给出具体的技能萃取工程方案。近期关于LLM智能体技能规范的研究开始关注技能的可移植性与安全性~\cite{skcc2025,skill_pseudocode2025,skill_comprehension2025}，但仍缺乏针对多智能体团队的完整技能生命周期管理方案。

本文提出的AgentsGroup2026技能萃取与赋予系统，正是对这一空白的系统性填补。系统的核心贡献在于：

\begin{enumerate}[nosep]
\item 设计并实现了七阶段WebSocket驱动的实时技能萃取流水线，将非结构化的多智能体对话转化为结构化的技能定义；
\item 构建了事件驱动的向量化技能索引器与带衰减策略的检索器，支持意图感知查询与角色推荐；
\item 提出了自然选择驱动的技能演化引擎，通过证据收集、LLM改进与A/B测试实现技能的达尔文式优化；
\item 建立了三池分类体系（特有/通用/储备）与技能身份归一化机制，使技能在团队间可继承、可进化。
\end{enumerate}

% ═══════════════════════════════════════════
\section{系统架构}
% ═══════════════════════════════════════════

AgentsGroup2026技能管理系统采用三层架构设计，如图\ref{fig:architecture}所示。三层分别对应技能生命周期的三个核心阶段——萃取、索引与演化——各层通过领域事件总线进行解耦通信。

\subsection{整体架构}

\textbf{第一层：技能萃取层（Extraction Layer）。}该层以技能萃取流水线（skill\_extract\_ws.py）为核心，通过WebSocket实时通信驱动七阶段萃取流程。智能体对话会话的原始内容被逐步转化为带有人工审核标记的结构化技能定义。萃取层是技能生命周期的入口——所有技能均诞生于真实的任务交互。

\textbf{第二层：索引与检索层（Indexing and Querying Layer）。}该层包含技能索引器（skill\_indexer.py）与衰减加权检索器（skill\_querier.py）。索引器订阅领域事件总线的SKILL\_CREATED/UPDATED/DELETED事件，采用TF-IDF向量化~\cite{salton1975vector}构建可搜索的技能向量空间。检索器在相似度得分基础上叠加时间衰减因子与使用频率增强，并提供意图映射与角色推荐功能。

\textbf{第三层：演化与赋予层（Evolution and Assignment Layer）。}该层包含技能演化引擎（skill\_evolver.py）、技能效果追踪器（skill\_tracker.py）和技能分类器（skill\_classifier.py）。演化引擎通过"证据→LLM改进→版本递增"的闭环实现技能指令的达尔文式优化。追踪器订阅TASK\_COMPLETED事件更新使用指标，自动将低效果技能降级。分类器基于演练证据将技能归入特有/通用/储备三池。演化层与AgentTeam模型的技能赋予机制集成，每个智能体携带技能列表，团队维护技能库。

\begin{figure}[h]
\centering
\fbox{\parbox{\linewidth}{\centering
\textbf{三层技能管理系统架构}\\[4pt]
\begin{tabular}{c}
\hline
\textbf{第三层：演化与赋予层} \\
技能演化引擎 | 效果追踪器 | 三池分类器 | 团队技能赋予 \\
\hline
\textbf{第二层：索引与检索层} \\
TF-IDF向量索引器 | 衰减加权检索器 | 意图感知查询 \\
\hline
\textbf{第一层：技能萃取层} \\
七阶段WebSocket流水线 | 时间轴UI推送 | 人工审核卡 \\
\hline
\end{tabular}
}}
\caption{AgentsGroup2026技能管理系统的三层架构。各层通过领域事件总线实现解耦通信。}
\label{fig:architecture}
\end{figure}

\subsection{事件驱动的解耦设计}

系统的三个层次通过统一的领域事件总线（EventBus）进行通信，这是实现层间解耦的关键设计。当技能萃取流水线完成一项萃取并生成新的技能定义后，索引器通过订阅SKILL\_CREATED事件自动将新技能纳入向量索引，无需萃取层主动调用索引接口。类似地，当智能体完成或失败任务时，追踪器通过订阅TASK\_COMPLETED/TASK\_FAILED事件更新技能的效果指标，进而触发演化引擎的建议生成。

这种事件驱动架构使得技能管理的各环节可以独立演进而互不阻塞——萃取层无需等待索引完成，演化层无需关注萃取细节。EventBus基于发布-订阅模式实现异步解耦，支持技能的完整生命周期追踪。

% ═══════════════════════════════════════════
\section{技能萃取流水线}
% ═══════════════════════════════════════════

技能萃取流水线是系统的入口引擎，负责将非结构化的多智能体对话会话转化为结构化技能定义。前端时间轴UI通过WebSocket与后端实时通信，用户可观察萃取过程的每一阶段进展。

\subsection{七阶段萃取流程}

萃取流水线按严格顺序执行七个阶段，如表\ref{tab:phases}所示。每个阶段由TimelinePhase枚举定义，阶段间通过锁定机制（lock\_reason）实现依赖管理。

\begin{table}[h]
\centering
\caption{七阶段萃取流水线}
\label{tab:phases}
\setlength{\tabcolsep}{4pt}
\begin{tabular}{cll}
\toprule
顺序 & 阶段 & 功能描述 \\
\midrule
1 & EXTRACT（抽取） & 从对话会话中提取候选技能片段 \\
2 & LLM\_PREFILL（LLM预填） & LLM生成技能名称、描述、指令初稿 \\
3 & REVIEW（审核） & 人类审核者评估LLM输出质量 \\
4 & APPROVE（批准） & 审核者批准或驳回技能定义 \\
5 & INDEX（索引） & 将批准的技能写入向量索引 \\
6 & VERIFY（验证） & 验证索引后的技能可被检索到 \\
7 & COMPLETE（完成） & 触发下游事件并推送给前端 \\
\bottomrule
\end{tabular}
\end{table}

每个阶段对应时间轴上的一个步骤节点。用户可看到每个步骤的进度百分比和锁状态（not\_started, in\_progress, waiting\_review, waiting\_approval, dependency\_blocked, completed）。锁定状态由前后阶段的依赖关系决定——例如，REVIEW阶段在LLM\_PREFILL完成前处于wait\_review状态；APPROVE阶段在审核完成前处于waiting\_approval状态。

\subsection{实时WebSocket推送机制}

萃取流水线的核心交互通过SkillExtractWSManager实现，该组件管理WebSocket连接池（按team\_id分组），提供四种推送消息类型：

\begin{itemize}[nosep]
\item \textbf{log\_node：}推送操作日志节点，包含阶段标签、状态、时间戳和元数据。前端在时间轴上渲染为可展开的步骤卡片。
\item \textbf{card\_update：}推送技能卡片的状态更新。卡片经历pending→recording→recorded→reviewing→reviewed→approved的生命周期，每次状态变更均触发实时推送。
\item \textbf{phase\_progress：}推送阶段步骤的进度百分比和锁状态。例如INDEX阶段在向量矩阵重建时显示"构建TF-IDF矩阵中..."的逐步进展。
\item \textbf{extraction\_completed：}萃取全部完成后推送汇总摘要，包含技能ID、名称、类别和索引状态。
\end{itemize}

WebSocket管理器采用单例模式，维护extraction\_id到萃取状态的缓存映射。客户端发送subscribe消息订阅特定萃取ID后，可回放已有的时间轴状态。这种设计确保了前端刷新或断连重连后状态的完整恢复。

\subsection{萃取状态的缓存与广播}

萃取状态以字典形式缓存于WebSocket管理器中。当LLM预填阶段生成技能指令初稿时，前端可通过时间轴UI实时看到LLM流式输出的文本逐步填充到卡片中。当人工审核者提交审核意见时，卡片状态从reviewing变为reviewed，同时推送review\_annotation内容。

一个典型的萃取会话可能同时处理多个候选技能——每个技能对应一条独立的时间轴路径（由node\_id区分）。前端以并行卡片组的形式展示，每条路径独立推进但共享同一个extraction\_id。完成卡片（completion card）包含realtime\_record（实时记录摘要）和review\_annotation（审核标注），形成完整的审计追踪。

% ═══════════════════════════════════════════
\section{技能索引与检索}
% ═══════════════════════════════════════════

\subsection{向量化索引}

技能索引器（SkillIndexer）采用事件驱动的异步索引架构。启动时，索引器通过EventBus订阅SKILL\_CREATED、SKILL\_UPDATED和SKILL\_DELETED三种事件，在事件发生时对内存中的向量索引进行增量更新。

索引器使用TF-IDF向量化~\cite{salton1975vector}将技能文本转为向量表示。每项技能文档由以下字段拼接而成：名称（权重×2以提升名称匹配的精确度）、描述、类别、slug标识符、所需工具列表、以及指令文本的前500字符。拼接文本通过sklearn的TfidfVectorizer进行特征提取，参数设置为max\_features=5000、ngram\_range=(1,2)、英文停用词过滤。

当sklearn不可用时，索引器自动降级为关键词匹配模式。降级模式下，检索器将查询词与文档文本进行分词匹配，以查询词在文档中的覆盖率作为相似度得分。这一设计使系统在轻量级部署环境下无需依赖完整的科学计算栈即可提供基本检索功能。

索引的向量矩阵采用延迟构建策略（lazy init）。当无文档更新时，已缓存的TF-IDF矩阵和skill\_ids列表保持不变；当发生文档增删改时，通过\_invalidate\_matrix标记矩阵失效，下次检索时自动触发重建。这种策略避免了每次更新后的全量重算，在高频更新的场景下显著降低了计算开销。

全量重建（rebuild）方法接受SkillStore记录的列表输入，清空现有文档并从头构建索引，适用于系统启动或数据迁移场景。

\subsection{衰减加权检索}

技能检索器（SkillQuerier）在索引器的原始相似度得分基础上叠加两层策略——时间衰减与频率增强——使查询结果的排序不仅依赖于文本相关性，还受到技能时效性与使用历史的调节。

检索器的核心得分公式为：

\begin{equation}
\mathrm{final\_score} = \mathrm{similarity} \times [\mathrm{decay} \times \omega_r + (1 - \omega_r)] \times \mathrm{frequency\_boost}
\label{eq:final_score}
\end{equation}

其中$\omega_r$为时效性权重（默认0.3），控制时间衰减对最终得分的影响程度。$\mathrm{decay}$为指数衰减因子：

\begin{equation}
\mathrm{decay} = e^{-\lambda \times \mathrm{age\_days}}
\label{eq:decay}
\end{equation}

其中$\lambda$为衰减系数（可配置，默认0.01），$\mathrm{age\_days}$为技能自上次更新以来的天数。最近更新的技能衰减较少，得分相对更高。

$\mathrm{frequency\_boost}$为频率增强因子：

\begin{equation}
\mathrm{frequency\_boost} = 1 + \alpha \times \log(1 + \mathrm{use\_count})
\label{eq:freq_boost}
\end{equation}

其中$\alpha$为频率增强系数（可配置，默认0.2），$\mathrm{use\_count}$为技能的累计使用次数。高频使用的技能获得对数级boost——这反映了"用得多的技能更可能是好技能"的启发式，但boost随使用次数增长的速度递减，避免已经过时的热门技能持续霸占检索结果。

检索器提供三级查询接口。基础query接口支持类别筛选、来源筛选、启用状态筛选等组合过滤器。query\_by\_intent接口将智能体的意图映射为优先技能类别——例如code意图优先检索general类技能，monitor意图优先检索automation类技能，research意图优先检索research类技能——当优先类别结果不足时自动补充通用结果。get\_recommendations接口基于智能体角色（developer/architect/researcher/qa/devops）生成推荐查询文本，并排除团队已有的技能以避免重复推荐。

检索器还维护各技能的使用统计（use\_count、last\_used\_at、first\_used\_at），提供record\_use和batch\_record\_use接口供上层（如任务执行器）在技能调用后更新统计数据，形成"使用→记录→增强"的正反馈循环。

\subsection{意图映射与角色推荐}

意图映射机制在query\_by\_intent方法中实现。查询意图（代码、研究、构建、分析、监控等）通过intent\_category\_map映射到技能类别：code→general（代码实现、测试、调试），research→research（需求分析、跨会话研究），monitor→automation（监控自动化），build→general（构建自动化），analyze→general（数据分析）。

角色推荐机制在get\_recommendations方法中实现。每种智能体角色预定义了领域查询词列表：developer角色查询"code implementation debugging refactoring testing"，architect角色查询"architecture design pattern interface definition"，qa角色查询"test design execution coverage regression"，devops角色查询"build automation container deployment monitoring"。角色查询结果经过团队已有技能的去重过滤，确保推荐的是团队中尚未覆盖的技能领域。

% ═══════════════════════════════════════════
\section{技能身份归一化}
% ═══════════════════════════════════════════

多智能体团队中，同一个技能可能以不同的名称或标识符出现在不同来源中——例如从萃取流水线生成的技能可能以"aws\_es\_scaling"为slug，而同一技能在熟练度存储中以"AWS ES Scaling"为名称，在技能库中以"es\_扩容"为别名。这种"同名不同ID"的问题会导致自然选择中的技能标识碎片化——同一个技能因ID不同而在多代演化中被视为不同基因，无法积累使用统计和演化历史。

AgentsGroup2026通过技能身份归一化机制解决这一问题。归一化过程包括：将所有技能名称转为小写、将空格和特殊分隔符（下划线、连字符、点号）统一替换、建立别名→规范ID的映射目录。当一个技能定义试图从非规范格式（如含大写字母、含特殊字符的昵称）引入时，身份归一化器自动将其匹配到已有的规范技能ID。

在实施中，AgentTeam模型的技能赋予接口在添加技能前先执行身份归一化——AgentProfile的skills列表在保存前经过canonicalize处理，确保所有技能引用指向规范ID。团队级技能库（team.skills字典）同样以规范ID为键，杜绝同一技能的多ID条目。这一机制从根本上防止了"多ID技能饥饿"——即排序靠前的技能ID因累积使用而获得evolutionary优势，而相同技能的其他ID因"未被使用"而被自然选择误判为低效。

% ═══════════════════════════════════════════
\section{技能演化引擎}
% ═══════════════════════════════════════════

\subsection{证据→归因→演化循环}

技能演化引擎（SkillEvolver）通过三个闭合步骤实现技能指令的达尔文式优化：（1）收集证据——从技能的使用计数、成功率、关联会话ID和用户反馈中提取演化依据；（2）LLM改进——将证据组装为提示词，调用LLM生成改进后的技能指令；（3）版本递增——将LLM输出作为新版本写回技能定义，版本号+1。

演化触发可通过两种模式：一是\textbf{人工触发}（manual），用户在界面上点击"演化"按钮；二是\textbf{自动建议}（auto\_low\_score），当技能追踪器检测到usage\_count≥5且effectiveness<0.4时自动生成改进建议。

LLM改进的效果取决于提示词的质量。系统提示词明确要求：保持核心意图、使指令更具体可操作、修正模糊表述、添加边界条件与异常处理指导、基于效果数据调整策略。当LLM不可用（chat\_harness未初始化或API调用失败）时，演化引擎降级为确定性改进——基于技能文本内容的关键词匹配，为特定领域技能（如RI成本优化）添加预设的改进步骤模板。

\subsection{技能效果追踪}

技能效果追踪器（SkillEffectivenessTracker）是演化引擎的度量基础设施。追踪器订阅EventBus的TASK\_COMPLETED和TASK\_FAILED事件，提取事件载荷中metadata.skills\_used列表，对每个被使用的技能更新四个核心指标：

\begin{itemize}[nosep]
\item \textbf{usage\_count}：技能被调用的总次数；
\item \textbf{success\_count / fail\_count}：成功/失败次数的累计；
\item \textbf{effectiveness} = success\_count / usage\_count，即时反映技能的实际效果；
\item \textbf{last\_used\_at}：最近使用时间戳。
\end{itemize}

追踪器内置自动降级机制：当usage\_count≥5且effectiveness<0.4时，技能的生命周期阶段（lifecycle\_stage）被自动置为DEGRADED，触发分类器将其移入储备池。每次任务完成事件还将session\_id追加到技能的evidence\_sessions列表中（保留最近50个），为后续的演化过程提供关联上下文。

\subsection{演化运行状态机}

在更高级的演化桥接器（EvolutionBridge）中，技能演化被建模为完整的演化运行（EvolutionRun），遵循七状态状态机：

\begin{equation}
\begin{aligned}
&\mathrm{IDENTIFYING} \rightarrow \mathrm{REFLECTING} \rightarrow \mathrm{MUTATING} \\
&\rightarrow \mathrm{AB\_TESTING} \rightarrow \mathrm{GATING} \rightarrow \mathrm{APPLIED}
\end{aligned}
\label{eq:evo_state}
\end{equation}

\textbf{IDENTIFYING阶段}通过聚合最近窗口（默认5个trial）的熟练度数据，识别成功率低于期望阈值（默认0.6）或趋势连续下滑的弱技能。\textbf{REFLECTING阶段}调用LLM对失败记录进行反思分析，生成结构化的反思报告。\textbf{MUTATING阶段}基于反思报告生成最多4个技能变体（candidates）。\textbf{AB\_TESTING阶段}在沙箱环境中对每个变体进行对照演练，收集适应度指标（fitness）和各维度得分。\textbf{GATING阶段}根据晋升条件判定获胜变体——适应度提升≥5\%且无单维度下降超过10\%。auto\_apply=False时停在GATING等待人工审批。

演化运行的触发模式支持三种：manual（手动触发）、auto\_low\_score（低分自动触发）、nightly（夜间定时触发）。全流程token预算由配置文件中的evolution\_budget\_tokens控制（默认200,000），防止演化过程消耗过量计算资源。

\subsection{技能合并与复制检测}

演化引擎提供了技能合并（merge）功能，用于消除技能库中的重复技能。合并策略支持两种模式：keep\_longest（保留指令文本最长的技能）和pick\_best\_score（保留效果最好的技能）。合并后的技能继承源技能的required\_tools并集、usage\_count和success\_count之和，effectiveness重新计算。

复制检测（find\_duplicates）基于相似度阈值（默认0.85）扫描技能库，识别指令文本高度相似但ID不同的技能对，生成合并建议。这为"同名不同ID"问题提供了补充的修复手段。

% ═══════════════════════════════════════════
\section{技能分类与赋予}
% ═══════════════════════════════════════════

\subsection{三池分类体系}

技能分类器（SkillClassifier）将团队的技能归入三个互斥的池：

\textbf{特有技能池（Exclusive Pool）}：单团队使用占比≥80\%、effectiveness≥0.6、且演练达标率满足rubric期望的技能。这类技能反映了团队的独特能力，是团队的核心竞争力。

\textbf{通用技能池（General Pool）}：至少2个团队采用或至少2个场景类目的演练达标、且发布门禁通过的技能。这类技能具有跨团队可迁移性，可被多个团队共享。

\textbf{储备技能池（Reserve Pool）}：新萃取尚未验证的技能默认为储备；effectiveness<0.4、超过90天未使用、或lifecycle\_stage=degraded的技能被强制回收至储备池。

分类判定遵循严格的优先级顺序：强制储备条件（degraded、低效果、过期、无使用记录）优先级最高；其次为通用条件（多团队采用或跨场景达标）；再次为特有条件（单团队高度依赖）。三个条件均不满足的技能留在储备池。

分类器带有基于滑动窗口的防抖机制：毕业（从储备→特有或从特有→通用）需要连续GRADUATE\_STREAK（默认2）个周期的即时分类达标；降级有DEMOTE\_GRACE（默认1）个周期的宽限，避免因单次数据波动导致的不稳定分类变更。

\subsection{技能赋予与团队集成}

技能赋予通过AgentTeam模型的skills字典实现。每个智能体（AgentProfile）携带skills列表，存储技能ID的引用列表。团队级技能库是团队内所有智能体共用的知识基础设施——智能体在执行任务时从自己的skills列表中加载技能指令作为系统提示的一部分。

技能赋予流程的关键步骤包括：（1）技能身份归一化，将候选技能的别名格式统一为规范ID；（2）去重检查，避免与智能体已有技能重复；（3）熟练度初始化，为新赋能的技能建立初次使用的熟练度记录。

团队技能库的分类视图（特有/通用/储备）为智能体队长（Coordinator）的任务分配决策提供依据——特有技能优先分配给持有该技能的智能体，通用技能可分配给任何具备匹配工具权限的成员，储备技能在被重分类至特有或通用池前不可用于主动分配。

% ═══════════════════════════════════════════
\section{技能可移植性分级}
% ═══════════════════════════════════════════

技能的可移植性表征其在不同的智能体运行环境中迁移部署的难度。技能注册中心（SkillRegistry）的classify\_portability方法对每个技能进行三级分类：

\textbf{一级（Tier 1）：纯提示词技能。}技能的指令文本不包含对任何外部工具或平台API的引用。这类技能可在任何LLM智能体框架间零成本迁移——只需将指令文本复制为系统提示。纯提示词技能是最高可移植性的技能形态，代表了跨平台技能复用的理想形式。

\textbf{二级（Tier 2）：CLI/API工具技能。}技能的指令文本中包含对工具调用的引用（如run\_python、run\_shell、web\_search、pip install），或者技能的required\_tools列表非空。这类技能迁移时需要目标环境提供等效工具集，但仍可能在工具接口标准化的情况下实现跨框架迁移。

\textbf{三级（Tier 3）：平台原生技能。}指令中包含对特定平台原语的引用（如channel.、marine\_base、process\_event、digital\_twin）。这类技能深度绑定特定平台的运行时环境，迁移成本极高，通常仅在平台内复用。

分类通过关键词匹配完成——平台标记优先于CLI标记，CLI标记优先于纯提示词。可移植性分级为技能在团队间的共享决策提供了工程依据：Tier 1技能可直接跨团队广播，Tier 2技能需附带工具依赖声明，Tier 3技能仅在兼容平台团队间共享。

% ═══════════════════════════════════════════
\section{实验分析}
% ═══════════════════════════════════════════

实验在AgentsGroup2026平台的完整技能生命周期中运行，覆盖AWS运维与构建系统两个任务领域、9次演化运行、45个智能体。实验目标为验证以下假设：

\begin{enumerate}[nosep]
\item 七阶段萃取流水线能从真实的多智能体对话会话中抽取结构化技能定义；
\item 向量索引与衰减加权检索能提供准确的技能匹配；
\item 演化引擎能通过LLM改进和A/B测试提升技能效能；
\item 三池分类体系能有效区分团队特有技能与跨团队通用技能。
\end{enumerate}

\subsection{萃取吞吐量}

七阶段萃取流水线的实时WebSocket推送机制使得前端时间轴UI能以亚秒级延迟展示各阶段的进展。各阶段耗时分布如下：EXTRACT阶段（0.3--0.8秒，取决于会话长度）主要消耗在对话文本的语义分割；LLM\_PREFILL阶段（1.2--3.5秒）取决于LLM推理延迟和候选技能数量；REVIEW和APPROVE阶段（人工操作，0.5--5分钟）是人机协作的瓶颈；INDEX和VERIFY阶段（0.1--0.3秒）在内存向量索引中瞬时完成。整体萃取一个会话（包含3--5个候选技能）的端到端耗时约为4--8分钟，其中人工审核占主导。

\subsection{索引延迟与检索精度}

TF-IDF索引的增量更新延迟在10--50毫秒量级（单文档更新触发矩阵失效→下次检索时重建）。全量重建的延迟与文档数量线性相关：100个技能全量重建约80毫秒，500个技能约350毫秒。由于矩阵采用延迟构建策略，高频更新场景下（如批量导入）的实际延迟低于每次更新即触发的全量重算。

检索精度方面，衰减加权策略在"近期更新技能优先"的测试场景中表现出预期效果：将recency\_weight设为0.3时，最近更新的技能平均排名提升2.1位。频率增强策略对"高频使用技能"的boost效果随使用次数对数增长——使用10次的技能获得约1.5倍的boost，使用100次的技能获得约1.9倍的boost，增幅递减防止了过时热门技能的垄断。

\subsection{演化效能}

演化引擎在9次演化运行中处理了45个智能体的技能演化建议。suggest\_evolution方法识别的改进候选与人工标注的一致性达78\%（30个改进建议中23个被人类审核员接受）。LLM改进后的技能指令在后续任务执行中的平均success\_rate提升为12\%（从67\%到75\%，基于建议被应用后的下一个10次使用统计），但样本量有限——23个被接受的改进中仅15个积累了足够的后优化使用数据。

A/B测试（AB\_TESTING阶段）中，变体技能与基线技能在相同沙箱场景下的对照演练表明，LLM生成的变体在37\%的案例中（11/30个变体对）表现出统计显著的适应度提升（fitness\_delta > 5\%）。

\subsection{技能继承模式}

三池分类体系中，9代混合竞争产生了5个dominant技能，其中包括跨域通用技能（interface\_definition）和领域特有技能（aws\_es\_scaling）。分类器正确地将interface\_definition归入通用池（2个团队采用，2个场景类目达标），将aws\_es\_scaling归入特有池（单团队使用占比92\%，effectiveness=0.73）。储备池中的技能在未积累足够使用证据前（usage\_count<5或effectiveness<0.4）不会被错误晋升，验证了防抖机制的稳定性。

可移植性分级的结果显示，5个dominant技能中2个为Tier 1（纯提示词），2个为Tier 2（需要工具支持），1个为Tier 3（绑定特定平台），这与AWS运维和构建系统的实际工具依赖分布一致。

% ═══════════════════════════════════════════
\section{相关工作}
% ═══════════════════════════════════════════

LLM驱动的多智能体系统在近年来取得了长足发展。AutoGen~\cite{wu2023autogen}提出基于对话编排的多智能体框架，允许开发者通过自然语言和代码灵活定义智能体交互模式。MetaGPT~\cite{hong2023metagpt}将人类软件工程中的标准操作流程（SOP）编码为提示词序列，使多智能体协作遵循结构化的角色分工。Generative Agents~\cite{park2023generative}构建了25个模拟人类日常行为的生成式智能体，展示了社会涌现行为的可能性。这些框架确立了多智能体协作的基础设施，但均未涉及从智能体交互中自动萃取技能的系统方法。

在技能学习方面，Voyager~\cite{wang2023voyager}的不断增长的代码技能库是该方向的开创性工作，但其技能以可执行代码形式存储，缺乏跨智能体、跨团队的泛化机制。Reflexion~\cite{shinn2023reflexion}通过语言反馈增强智能体的学习能力，但反思结果未结构化为可复用的技能知识库。ReAct~\cite{yao2022react}探索了推理与行动的交织生成，其推理追踪可视为技能的雏形但未被显式建模。

在技能可移植性方面，近期工作开始系统化地审视LLM智能体技能规范的工程问题。SkCC~\cite{skcc2025}引入编译器设计思想，将格式无关的技能Markdown编译为目标框架特定的提示词，首次解决了跨框架技能迁移的格式适配问题。Skill-as-Pseudocode~\cite{skill_pseudocode2025}观察到Markdown技能库导致智能体陷入"困惑→重检索→仍困惑"的循环，提出将技能转换为类型化伪代码以提高机器可读性。技能规范用户理解相关研究~\cite{skill_comprehension2025}揭示了当前技能规范在输出契约、边界披露、示例能力演示上的信息缺失问题。本系统的可移植性分级（Tier 1--3）从另一个维度——运行时依赖——补充了这些工作。

CoALA~\cite{sumers2023coala}从认知科学和符号人工智能的丰富历史出发，提出了语言智能体的统一认知架构，将技能定位为长期记忆中的程序性知识。本系统可视为CoALA框架的技能管理模块的一个具体工程实现——从会话萃取（感知）、向量索引（记忆存储）、衰减检索（记忆检索）到演化赋予（学习与改进）构成了完整的认知循环。

在演化计算方面，达尔文自然选择理论~\cite{darwin1859}及其现代扩展（扩展进化综合~\cite{pigliucci2010extended,laland2015extended}、建设性中性演化~\cite{stoltzfus1999constructive}）和数字进化实验（Avida~\cite{lenski2003evolutionary}）为本文的演化引擎提供了理论基础。Boyd与Richerson的双继承理论~\cite{boyd1985culture}——基因演化与文化演化的双重传递机制——映射为本文中技能基因的垂直遗传（跨代交叉与变异）与水平传递（团队间技能共享）。

% ═══════════════════════════════════════════
\section{结论}
% ═══════════════════════════════════════════

本文提出了面向多智能体团队的完整技能萃取与赋予系统，覆盖了从技能诞生（萃取流水线）、成长（向量索引与衰减检索）、成熟（演化引擎）到传承（三池分类与团队赋予）的全生命周期。系统以事件驱动架构和WebSocket实时通信为基础，实现了非结构化对话→结构化技能→可检索知识→可演化基因的完整闭环。

七阶段萃取流水线通过WebSocket驱动的时间轴UI将人工审核嵌入萃取流程，在自动化与人工质量把关之间取得了平衡。向量索引与衰减加权检索器为技能的"可发现性"提供了工程基础——TF-IDF负责文本匹配，时间衰减负责时效排序，频率增强负责使用反馈。演化引擎将自然选择的逻辑映射到技能指令优化中——证据收集是"环境观察"，LLM改进是"变异生成"，A/B测试是"环境选择"，而版本递增是"性状遗传"。

三池分类体系与可移植性分级为技能的"可迁移性"提供了量化框架。特有/通用/储备的分类回答了"这个技能属于谁"的问题；Tier 1--3的分级回答了"这个技能能去哪里"的问题。两者共同构成了多智能体团队技能治理的基础设施。

实验结果表明，系统在萃取吞吐量、检索精度、演化效能和分类准确性方面均达到预期水平。9代混合竞争产生的dominant技能列表验证了技能跨代遗传的可行性。未来的工作将扩展至更大规模的智能体种群（≥100）、更多样的任务领域（≥5）、以及更长时间的演化窗口（≥20代），以更全面地评估技能萃取与演化系统的泛化性与可扩展性。

\bibliographystyle{IEEEtran}
\bibliography{skill_extraction_references}

\end{document}
